% Options for packages loaded elsewhere
\PassOptionsToPackage{unicode}{hyperref}
\PassOptionsToPackage{hyphens}{url}
\PassOptionsToPackage{dvipsnames,svgnames,x11names}{xcolor}
\documentclass[
]{article}
\usepackage{xcolor}
\usepackage[margin=1in]{geometry}
\usepackage{amsmath,amssymb}
\setcounter{secnumdepth}{5}
\usepackage{iftex}
\ifPDFTeX
  \usepackage[T1]{fontenc}
  \usepackage[utf8]{inputenc}
  \usepackage{textcomp} % provide euro and other symbols
\else % if luatex or xetex
  \usepackage{unicode-math} % this also loads fontspec
  \defaultfontfeatures{Scale=MatchLowercase}
  \defaultfontfeatures[\rmfamily]{Ligatures=TeX,Scale=1}
\fi
\usepackage{lmodern}
\ifPDFTeX\else
  % xetex/luatex font selection
\fi
% Use upquote if available, for straight quotes in verbatim environments
\IfFileExists{upquote.sty}{\usepackage{upquote}}{}
\IfFileExists{microtype.sty}{% use microtype if available
  \usepackage[]{microtype}
  \UseMicrotypeSet[protrusion]{basicmath} % disable protrusion for tt fonts
}{}
\makeatletter
\@ifundefined{KOMAClassName}{% if non-KOMA class
  \IfFileExists{parskip.sty}{%
    \usepackage{parskip}
  }{% else
    \setlength{\parindent}{0pt}
    \setlength{\parskip}{6pt plus 2pt minus 1pt}}
}{% if KOMA class
  \KOMAoptions{parskip=half}}
\makeatother
\setlength{\emergencystretch}{3em} % prevent overfull lines
\providecommand{\tightlist}{%
  \setlength{\itemsep}{0pt}\setlength{\parskip}{0pt}}
\usepackage[dvipsnames]{xcolor}
\usepackage{booktabs}
\usepackage{longtable}
\usepackage{array}
\usepackage{microtype}
\usepackage{bookmark}
\IfFileExists{xurl.sty}{\usepackage{xurl}}{} % add URL line breaks if available
\urlstyle{same}
\hypersetup{
  pdftitle={Neural Field Beta},
  pdfauthor={Ben Tang \textbar{} Greg Cogan Lab, Duke University},
  colorlinks=true,
  linkcolor={NavyBlue},
  filecolor={Maroon},
  citecolor={Blue},
  urlcolor={NavyBlue},
  pdfcreator={LaTeX via pandoc}}

\title{Neural Field Beta}
\usepackage{etoolbox}
\makeatletter
\providecommand{\subtitle}[1]{% add subtitle to \maketitle
  \apptocmd{\@title}{\par {\large #1 \par}}{}{}
}
\makeatother
\subtitle{Minimum Viable Sequential Test of Coordinate-Aware
Cross-Patient Transfer}
\author{Ben Tang \textbar{} Greg Cogan Lab, Duke University}
\date{April 2026}

\begin{document}
\maketitle

{
\setcounter{tocdepth}{3}
\tableofcontents
}
\section{Status}\label{status}

This document is the \textbf{beta} version of the neural-field
direction. Its purpose is not to realize the full Neural Field Perceiver
at once. Its purpose is to test the smallest defensible version of the
hypothesis:

\begin{quote}
\textbf{Minimum viable hypothesis.} Coarse anatomical coordinates
improve cross-patient LOPO speech decoding when added to the current
best-performing local spatial model, under the existing grouped-by-token
evaluation protocol.
\end{quote}

The design is deliberately conservative. It preserves the current
strongest LOPO recipe wherever possible and changes one thing at a time.

\section{Why A Beta Document Exists}\label{why-a-beta-document-exists}

The full \texttt{v10} design is much cleaner than \texttt{v9}, but it
still bundles several unresolved bets:

\begin{itemize}
\tightlist
\item
  It introduces a new encoder family, new latent geometry, virtual
  electrodes, attention-based readout, geometric correction, and two
  reconstruction losses in one proposal.
\item
  It still relies on raw volumetric MNI geometry as the main spatial
  scaffold, even though the broader project notes argue that MNI is
  trustworthy mainly at coarse cross-patient scale, while local
  within-patient structure should still come from the grid.
\item
  Its cross-patient reconstruction loss remains a weak and ambiguous
  supervisory signal: same-token pairing plus full-trial averaging can
  regularize, but it is not a clean first test of whether coordinates
  help.
\end{itemize}

The beta design removes those confounds. If the coordinate hypothesis is
real, the beta path should show it with far less architectural risk.

\section{Design Principles}\label{design-principles}

\begin{enumerate}
\def\labelenumi{\arabic{enumi}.}
\tightlist
\item
  Preserve the current LOPO winner unless a change is required to test
  the hypothesis.
\item
  Use coordinates only at the scale where they are scientifically
  defensible.
\item
  Keep local within-patient topology in the model.
\item
  Add no self-supervised or cross-patient reconstruction loss until
  plain supervised coordinate conditioning shows signal.
\item
  Gate every additional mechanism on a prior positive result.
\item
  Judge success on \textbf{all patients}, not on S14 alone.
\end{enumerate}

\section{Non-Goals}\label{non-goals}

This beta architecture explicitly does \textbf{not} try to test all of
the following at once:

\begin{itemize}
\tightlist
\item
  Perceiver-style latent arrays
\item
  Virtual electrodes
\item
  Cross-patient reprojection
\item
  Trial-matched latent swapping
\item
  Patient embeddings in the encoder
\item
  Point-cloud attention as the first implementation
\item
  Full neural-field identifiability claims
\end{itemize}

If the beta path fails, that is evidence against the immediate value of
the coordinate hypothesis in this data regime. It is not evidence that a
larger transformer would have rescued it.

\section{Anchor Model}\label{anchor-model}

The anchor is the \textbf{current best LOPO pipeline}, unchanged except
where the beta stages explicitly intervene:

\begin{itemize}
\tightlist
\item
  Existing grouped-by-token LOPO protocol
\item
  Existing multi-scale temporal frontend
\item
  Existing best local spatial read-in / pooling path
\item
  Existing articulatory bottleneck classification head
\item
  Existing LP-FT target adaptation regime
\item
  Existing evaluation recipe and reporting contract
\end{itemize}

This matters. The beta comparison must answer:

\begin{quote}
Does adding coarse anatomical information help \textbf{the model that
already works best}, rather than a brand-new model family?
\end{quote}

\section{Core Architectural Decision}\label{core-architectural-decision}

The beta architecture uses \textbf{dual spatial encoding}:

\begin{itemize}
\tightlist
\item
  \textbf{Local spatial structure} comes from the patient's native grid
  topology, as in the current pipeline.
\item
  \textbf{Cross-patient spatial structure} comes from coarse anatomical
  coordinates attached to each electrode.
\end{itemize}

This matches the project's theoretical framing:

\begin{itemize}
\tightlist
\item
  MNI is useful at coarse somatotopic scale.
\item
  Grid adjacency is useful at fine within-array scale.
\item
  Neither should be forced to do the other's job.
\end{itemize}

\section{Beta Architecture}\label{beta-architecture}

\subsection{1. Inputs}\label{inputs}

For patient \(i\), each trial provides:

\begin{itemize}
\tightlist
\item
  HGA grid \(X^{(i)} \in \mathbb{R}^{H_i \times W_i \times T}\)
\item
  Electrode-validity mask \(M^{(i)} \in \{0,1\}^{H_i \times W_i}\)
\item
  Static coordinate maps
  \(C^{(i)} \in \mathbb{R}^{H_i \times W_i \times c_{\text{coord}}}\)
\end{itemize}

The coordinate channels are attached only at valid electrode locations.
Dead grid positions remain zeroed.

\subsection{2. Coordinate
Representation}\label{coordinate-representation}

The beta path uses \textbf{coarse coordinates only}.

Stage-dependent coordinate variants:

\begin{itemize}
\tightlist
\item
  \texttt{Beta-1}: standardized raw \((x, y, z)\) MNI channels
\item
  \texttt{Beta-2+}: coarse Fourier features with \(F = 3\), not
  \(F = 4\)
\end{itemize}

Rationale:

\begin{itemize}
\tightlist
\item
  Raw XYZ is the simplest possible coordinate test.
\item
  \(F=3\) allows smooth coarse nonlinear spatial structure.
\item
  The finest \texttt{v10} Fourier scale is intentionally excluded from
  the beta path because it pushes below plausible cross-patient
  correspondence precision.
\end{itemize}

\subsection{3. Fusion Strategy}\label{fusion-strategy}

The model remains grid-native. Coordinates are \textbf{added as static
side channels}, not as a replacement for the grid model.

For each trial:

\begin{enumerate}
\def\labelenumi{\arabic{enumi}.}
\item
  Broadcast coordinate channels across time.
\item
  Concatenate activity and coordinates:

  \[
  U^{(i)} = \mathrm{concat}\bigl(X^{(i)}, C^{(i)}_{\text{broadcast}}\bigr)
  \]
\item
  Apply a shallow shared channel mixer:

  \[
  \tilde{U}^{(i)} = \mathrm{Conv2d}_{1 \times 1}(U^{(i)})
  \]
\item
  Feed the mixed tensor into the \textbf{existing} local spatial read-in
  path.
\end{enumerate}

Interpretation:

\begin{itemize}
\tightlist
\item
  The current model still learns local spatial filters on the patient's
  array.
\item
  Coordinates tell the read-in what cortical zone each local pattern
  likely belongs to.
\item
  The architecture never asks MNI to substitute for local grid
  structure.
\end{itemize}

\subsection{4. Backbone and Decoder}\label{backbone-and-decoder}

Backbone and phoneme decoder remain unchanged from the current LOPO
winner.

This is a hard rule of the beta design:

\begin{itemize}
\tightlist
\item
  no new temporal backbone
\item
  no new latent bottleneck
\item
  no new query readout
\item
  no new inference recipe
\end{itemize}

If coordinates help, they should help \textbf{before} a backbone
rewrite.

\section{Sequential Experiment
Ladder}\label{sequential-experiment-ladder}

\subsection{Beta-0: Population Anchor}\label{beta-0-population-anchor}

\textbf{Question:} What is the true all-patient anchor under the current
best LOPO recipe?

Model:

\begin{itemize}
\tightlist
\item
  Current best LOPO architecture
\item
  No coordinates
\item
  No new losses
\end{itemize}

Purpose:

\begin{itemize}
\tightlist
\item
  Establish the population baseline
\item
  Quantify patient-wise variability
\item
  Prevent S14-specific overinterpretation
\end{itemize}

Gate:

\begin{itemize}
\tightlist
\item
  Required before any architectural conclusion
\end{itemize}

\subsection{Beta-1: Minimal Coordinate
Injection}\label{beta-1-minimal-coordinate-injection}

\textbf{Question:} Do coarse anatomical coordinates help at all?

Model:

\begin{itemize}
\tightlist
\item
  Anchor model
\item
  Add standardized raw XYZ channels
\item
  Shared \texttt{1x1} fusion layer
\item
  No \(\Delta\), no \(\omega\), no reconstruction loss
\end{itemize}

This is the minimum viable test.

Why this first:

\begin{itemize}
\tightlist
\item
  It isolates the value of coordinates without importing geometric
  correction, latent inversion, or auxiliary supervision.
\item
  It preserves the local spatial inductive bias already validated by the
  current model.
\item
  Failure here is highly informative: if raw coarse coordinates do not
  help the current best model, the case for a much larger
  coordinate-aware stack weakens substantially.
\end{itemize}

Primary comparison:

\begin{itemize}
\tightlist
\item
  \texttt{Beta-1} vs \texttt{Beta-0}
\end{itemize}

Success gate:

\begin{itemize}
\tightlist
\item
  Mean PER improves by at least \texttt{3pp}
\item
  Effect sign is positive for a majority of patients
\item
  Paired patient-wise test passes the pre-registered threshold
\end{itemize}

Stop rule:

\begin{itemize}
\tightlist
\item
  If \texttt{Beta-1} fails, stop the coordinate path and prioritize data
  pooling / evaluation fixes instead
\end{itemize}

\subsection{Beta-2: Lightweight Geometric
Correction}\label{beta-2-lightweight-geometric-correction}

\textbf{Question:} Is residual rigid misregistration the missing piece
after plain coordinates?

Model:

\begin{itemize}
\item
  \texttt{Beta-1}
\item
  Add per-patient rigid correction on coordinates only:

  \[
  \tilde{\mathbf{p}}_j^{(i)} = R(\omega^{(i)})\,\mathbf{p}_j^{(i)} + \Delta^{(i)}
  \]
\item
  Use strong regularization toward identity:

  \[
  \mathcal{L}_{\text{geom}} = \lambda_{\Delta}\|\Delta^{(i)}\|_2^2 + \lambda_{\omega}\|\omega^{(i)}\|_2^2
  \]
\end{itemize}

Rules:

\begin{itemize}
\tightlist
\item
  Coordinates only are transformed
\item
  No patient embedding
\item
  No decoder-side patient code
\end{itemize}

Why this second:

\begin{itemize}
\tightlist
\item
  It directly tests whether the coordinate problem is mostly ``use
  anatomy'' or ``correct anatomy.''
\item
  It is still low-dimensional and interpretable.
\end{itemize}

Primary comparison:

\begin{itemize}
\tightlist
\item
  \texttt{Beta-2} vs \texttt{Beta-1}
\end{itemize}

Success gate:

\begin{itemize}
\tightlist
\item
  Consistent positive patient-wise effect
\item
  Learned transforms remain small and anatomically plausible
\end{itemize}

Stop rule:

\begin{itemize}
\tightlist
\item
  If transforms grow large or unstable, treat them as a failed
  identifiability channel and remove them
\end{itemize}

\subsection{Beta-3: Coarse Nonlinear Coordinate
Encoding}\label{beta-3-coarse-nonlinear-coordinate-encoding}

\textbf{Question:} Are raw XYZ channels too linear to expose the
coordinate signal?

Model:

\begin{itemize}
\tightlist
\item
  Best of \texttt{Beta-1} / \texttt{Beta-2}
\item
  Replace raw XYZ with coarse Fourier features \(F=3\)
\item
  Keep the rest unchanged
\end{itemize}

Why this is third:

\begin{itemize}
\tightlist
\item
  It tests nonlinear spatial encoding only after raw coordinates have
  shown some value
\item
  It avoids blaming Fourier engineering for failure of the underlying
  hypothesis
\end{itemize}

Primary comparison:

\begin{itemize}
\tightlist
\item
  \texttt{Beta-3} vs best prior beta stage
\end{itemize}

Success gate:

\begin{itemize}
\tightlist
\item
  Small but consistent incremental benefit
\item
  No evidence that performance gain comes only from one or two patients
\end{itemize}

\subsection{Beta-4: Within-Patient Masked Reconstruction
Only}\label{beta-4-within-patient-masked-reconstruction-only}

\textbf{Question:} Does a light reconstruction objective improve spatial
robustness once supervised coordinate conditioning already works?

Model:

\begin{itemize}
\tightlist
\item
  Best of \texttt{Beta-1} through \texttt{Beta-3}
\item
  Add a \textbf{within-patient only} masked-electrode reconstruction
  head
\item
  Reconstruct held-out electrode activity from the same patient's
  remaining electrodes
\end{itemize}

Loss:

\[
\mathcal{L} = \mathcal{L}_{\text{class}} + \lambda_{\text{recon}} \mathcal{L}_{\text{masked-recon}}
\]

Properties:

\begin{itemize}
\tightlist
\item
  No cross-patient trial pairing
\item
  No token matching
\item
  No patient latent swapping
\item
  No full neural-field decoder
\end{itemize}

Why this is the first auxiliary loss:

\begin{itemize}
\tightlist
\item
  It is operationally simple
\item
  It tests spatial interpolation robustness without creating a new
  ambiguous cross-patient objective
\item
  It uses trustworthy within-patient structure before touching noisy
  cross-patient reconstruction
\end{itemize}

Primary comparison:

\begin{itemize}
\tightlist
\item
  \texttt{Beta-4} vs best prior beta stage
\end{itemize}

Stop rule:

\begin{itemize}
\tightlist
\item
  If reconstruction helps the auxiliary metric but not PER, remove it
\end{itemize}

\subsection{Beta-5: Escalation
Threshold}\label{beta-5-escalation-threshold}

Only after the above sequence should the project consider:

\begin{itemize}
\tightlist
\item
  virtual electrodes
\item
  point-set / Perceiver attention
\item
  cross-patient reprojection
\item
  patient-aware decoder terms
\item
  trial-swapped latent supervision
\end{itemize}

The escalation condition is:

\begin{quote}
Coordinates show clear population-level value in the beta path, but the
remaining gap to the anchor still appears spatial rather than purely
data-limited.
\end{quote}

\section{Why Cross-Patient Reprojection Is Excluded From
Beta}\label{why-cross-patient-reprojection-is-excluded-from-beta}

Cross-patient reprojection is not the minimum viable test because it
adds several ambiguities at once:

\begin{itemize}
\tightlist
\item
  same-token matching can reward token templates
\item
  trial-level averaging discards temporal structure
\item
  cross-patient residuals are large and hard to interpret
\item
  failure is ambiguous: alignment problem, loss-noise problem, or
  architecture problem
\end{itemize}

That loss may still be worth exploring later. It is not the right first
experiment.

\section{Training Protocol}\label{training-protocol}

\subsection{Fixed Across Beta-0 Through
Beta-4}\label{fixed-across-beta-0-through-beta-4}

\begin{itemize}
\tightlist
\item
  Same grouped-by-token splits
\item
  Same LOPO patient reporting
\item
  Same target adaptation regime
\item
  Same seeds
\item
  Same augmentation policy unless an ablation explicitly changes it
\item
  Same evaluation recipe
\end{itemize}

This is necessary because the current evidence shows recipe variance is
already large enough to obscure small architecture effects.

\subsection{Optimization Policy}\label{optimization-policy}

\begin{itemize}
\tightlist
\item
  Reuse the current optimizer and scheduler unless a stage introduces a
  new parameter type
\item
  New geometric parameters use lower LR than backbone weights
\item
  Coordinate channels themselves are deterministic inputs, not learned
  embeddings, until \texttt{Beta-3}
\end{itemize}

\section{Evaluation Contract}\label{evaluation-contract}

Primary metric:

\begin{itemize}
\tightlist
\item
  Mean patient-wise PER under LOPO LP-FT
\end{itemize}

Required reporting:

\begin{itemize}
\tightlist
\item
  Per-patient PER table
\item
  Mean and median PER
\item
  Seed variance
\item
  Paired comparisons against the immediately previous beta stage
\end{itemize}

Secondary diagnostics:

\begin{itemize}
\tightlist
\item
  Whether gains concentrate in low-coverage or high-coverage patients
\item
  Whether gains are larger for displaced arrays
\item
  Whether rigid correction magnitude correlates with known registration
  uncertainty
\end{itemize}

\section{Decision Logic}\label{decision-logic}

\subsection{If Beta-1 Fails}\label{if-beta-1-fails}

Interpretation:

\begin{itemize}
\tightlist
\item
  Coarse coordinates do not add enough value to the current best local
  model in this regime
\end{itemize}

Action:

\begin{itemize}
\tightlist
\item
  Stop architecture escalation
\item
  Prioritize cross-task pooling and population evaluation
\end{itemize}

\subsection{If Beta-1 Works but Beta-2/Beta-3
Fail}\label{if-beta-1-works-but-beta-2beta-3-fail}

Interpretation:

\begin{itemize}
\tightlist
\item
  Coordinates help, but only in a simple coarse form
\end{itemize}

Action:

\begin{itemize}
\tightlist
\item
  Keep the simple coordinate-conditioned model
\item
  Do not escalate to the full Neural Field Perceiver
\end{itemize}

\subsection{If Beta-1 Through Beta-4 All
Help}\label{if-beta-1-through-beta-4-all-help}

Interpretation:

\begin{itemize}
\tightlist
\item
  The coordinate hypothesis is real and robust enough to justify larger
  architectural investment
\end{itemize}

Action:

\begin{itemize}
\tightlist
\item
  Re-open the Perceiver / latent-field direction with much stronger
  prior evidence
\end{itemize}

\section{Recommended Initial Build}\label{recommended-initial-build}

If only one beta experiment is implemented first, it should be:

\subsection{\texorpdfstring{\texttt{Beta-1}}{Beta-1}}\label{beta-1}

\begin{itemize}
\tightlist
\item
  Current best LOPO model
\item
  Add standardized XYZ coordinate channels
\item
  Shared \texttt{1x1} coordinate-activity fusion
\item
  No geometric correction
\item
  No reconstruction loss
\item
  Full all-patient evaluation
\end{itemize}

This is the cleanest first test of the hypothesis and the one least
likely to produce uninterpretable failure.

\section{Summary}\label{summary}

The beta program asks one question at a time:

\begin{enumerate}
\def\labelenumi{\arabic{enumi}.}
\tightlist
\item
  Do coordinates help at all?
\item
  If yes, does rigid correction help?
\item
  If yes, does coarse nonlinear encoding help?
\item
  If yes, does a simple within-patient reconstruction loss help?
\item
  Only then: is a larger field model justified?
\end{enumerate}

That sequence is scientifically stronger than starting with the full
Neural Field Perceiver. It preserves the current best model, respects
the actual reliability scale of MNI coordinates, and maximizes what can
be learned from either a positive or a negative result.

\end{document}
